\section{Introduction}

Reed--Solomon proximity questions connect three quantitatively different
objects:
\begin{enumerate}[label=(\roman*)]
  \item the relative distance of a received word from a code;
  \item the number of codewords in a ball at that radius;
  \item the fraction of affine challenges that conceal distance of an
    underlying tuple of received words.
\end{enumerate}
The Grand List Decoding and Grand Mutual Correlated Agreement challenges ask
for sharp forms of the latter two quantities for constant-rate Reed--Solomon
codes over smooth evaluation domains \cite{arnonBonehFenzi2026}.

This paper treats the coding problem independently of any protocol or
application-specific source of received words.  Its universal statements
quantify over every word in the ambient alphabet and every affine line.  A
later application may map structured observations into this ambient space, but
that map is not part of the coding theorem.

\subsection{Central organization}

Let
\[
  \mathcal C=\RS[\F,\mathcal L,k]\subseteq\F^n,
  \qquad n=|\mathcal L|,
  \qquad \rho=\frac{k}{n}.
\]
The two-fold interleaving is
\[
  \Cpair\subseteq(\F^2)^n.
\]
Every pair word \(p\in(\F^2)^n\) determines an affine line
\[
  y_\gamma(p)=p_0+\gamma p_1\in\F^n,
  \qquad \gamma\in\F.
\]
The intended proof architecture is therefore
\[
\begin{array}{c}
  p\in(\F^2)^n \\[2mm]
  \text{distance and list size relative to }\Cpair
\end{array}
\quad\xrightarrow{\ \gamma\in\F\ }\quad
\begin{array}{c}
  y_\gamma(p)\in\F^n \\[2mm]
  \text{distance relative to }\mathcal C .
\end{array}
\]
List decoding controls ambiguity of the full pair; MCA controls the fraction
of challenges for which a scalar observation fails to expose pair-farness.

\subsection{Claim discipline}

The plane/pair/line identifications below are exact equivalences.  They ensure
that no received words or affine lines are lost by the chosen presentation.
They do not prove capacity-radius list decoding or MCA.  Those quantitative
statements are separately labeled as proof targets until a complete argument
and its parameter inequalities have been supplied.

In particular, a statement of the form
\[
  \dist(p,\Cpair)\ge\delta
\]
is not a list-size bound, and neither fact alone implies a security estimate
of \(2^{-128}\).  The required list size or bad-challenge count must be compared
explicitly with \(\epsstar|\F|\).

\subsection{Contributions sought}

The completed paper is intended to establish:
\begin{enumerate}[label=(\roman*)]
  \item a capacity-radius list-size estimate for
    \(\mathcal C^{\equiv m}\), with \(m=2\) as the primary plane endpoint;
  \item a mutual-correlated-agreement estimate at the same normalized radius;
  \item explicit field-size hypotheses sufficient to obtain a selected
    \(\epsstar\);
  \item a machine-checked formalization of the definitions, reductions, and
    final quantitative theorem.
\end{enumerate}
